\lettrine{E}{arly} childhood spanning the first five years of life represents a critical period of human brain development, during which rapid and tightly coordinated processes---including dendritic arborization, synapse proliferation, glial differentiation, and axonal myelination---reshape cortical circuits at multiple scales\cite{petanjekExtraordinaryNeotenySynaptic2011,xiaDevelopmentSensorimotorvisualConnectome2024,tooleyStructureFunctionCoupling2025}. Over this period, the cortex transitions from a relatively immature architecture to one organized into large-scale structural and functional systems, including canonical networks and gradients that span sensorimotor to higher-order association regions\cite{xuStructuralConnectivityMatures2024,babaeeghazviniBrainStructuralFunctional2021}. This macro-scale organization supports the emergence of core cognitive functions and is sculpted by intense experience-dependent plasticity, rendering infancy both a window of opportunity and a period of heightened neurodevelopmental risk\cite{zhangDynamicStructureFunction2024}. Disruptions to how structural wiring and functional interactions become aligned are thought to contribute to vulnerability for later neurological and psychiatric disorders, including stroke, Parkinson's disease, schizophrenia, epilepsy, bipolar disorder, multiple sclerosis, mild cognitive impairment, and attention-deficit/hyperactivity disorder\cite{gonzalez-ramirezBiologicallyConstrainedMathematical2015,parkConnectomewideStructurefunctionCoupling2023,somanCorticalStructuralFunctional2023,chiangStructuralFunctionalCoupling2015,sunModularlevelAlterationsStructure2017}. Establishing normative trajectories of cortex-wide structure-function coupling during this sensitive period is therefore essential for understanding how early brain architecture supports healthy cognition and for identifying deviations that may presage atypical development.

Magnetic resonance imaging (MRI) provides a noninvasive approach to characterize the developing brain by measuring signals from hydrogen nuclei in tissue water and lipids. In a strong static magnetic field (\(B_0\)), proton spins partially align with the field and precess at the Larmor frequency. A radiofrequency pulse perturbs this equilibrium, and the emitted signal during relaxation is sampled to form an image. Tissue contrast arises because gray matter, white matter, and cerebrospinal fluid differ in longitudinal (\(T_1\)) and transverse (\(T_2/T_2^*\)) relaxation behavior. Spatial localization is encoded through magnetic field gradients that map frequency and phase to position, allowing reconstruction from \(k\)-space by inverse Fourier transform. In practice, acquisition parameters such as repetition time, echo time, and flip angle determine how strongly each relaxation component contributes to image contrast. These principles underpin T1-weighted and T2-weighted imaging, which are central to infant neuroimaging pipelines for cortical surface reconstruction, tissue segmentation, and myelination-sensitive measurements. MRI is especially valuable in pediatric cohorts because it avoids ionizing radiation while preserving whole-brain coverage and high anatomical detail\cite{westbrookMRIPractice6th}.

Functional MRI (fMRI) extends the same MR physics to study brain function through blood-oxygen-level-dependent (BOLD) contrast. BOLD signal changes arise because oxyhemoglobin and deoxyhemoglobin differ in magnetic susceptibility, which modulates local \(T_2^*\)-weighted signal intensity. When local neuronal activity increases, neurovascular coupling produces changes in cerebral blood flow, blood volume, and oxygen metabolism; the resulting reduction in deoxyhemoglobin typically increases BOLD signal. Therefore, fMRI does not measure neuronal firing directly, but rather a hemodynamic correlate of neural activity with second-level temporal dynamics and millimeter-level spatial specificity\cite{westbrookMRIPractice6th}. In resting-state fMRI (rs-fMRI), low-frequency spontaneous BOLD fluctuations are acquired without an explicit task, and temporal correlations among regions are used to estimate functional connectivity. These correlation structures can be represented as discrete networks or continuous cortical gradients, providing complementary views of macroscale brain organization. In infancy, this framework is powerful for mapping how sensorimotor and association systems differentiate over time, but it is technically demanding because infant data are particularly vulnerable to motion artifacts, susceptibility distortions, and low signal-to-noise ratio. Rigorous preprocessing and denoising are therefore required to recover biologically meaningful trajectories from rs-fMRI, including motion correction, distortion correction, anatomical registration, nuisance removal, and strict quality control. Within this context, fMRI offers an essential systems-level readout for linking structural scaffolds, cortical dynamics, and emerging functional hierarchies during early postnatal development\cite{westbrookMRIPractice6th}.

%Structure, function, and structure–function coupling
At the macroscale, cortical communication can be described in terms of structural connectivity (SC), which reflects the architecture of white-matter pathways linking regions, and functional connectivity (FC), which captures the temporal coupling of spontaneous activity between those regions\cite{gaoDevelopmentHumanBrain2015,sylvesterNetworkspecificSelectivityFunctional2023,somanCorticalStructuralFunctional2023}. During childhood and adolescence, FC patterns become increasingly modular and segregated, differentiating sensorimotor, attentional, and association systems, while SC undergoes parallel refinement with strengthening and integration of long-range tracts\cite{sylvesterNetworkspecificSelectivityFunctional2023,sydnorNeurodevelopmentAssociationCortices2021}. The degree to which regional patterns of FC follow the underlying SC—termed structure-function coupling (SFC)—provides a compact index of how efficiently anatomical wiring supports functional interactions\cite{baumDevelopmentStructureFunction2020,sylvesterNetworkspecificSelectivityFunctional2023}. In adults, SFC is not uniform across the cortex: early-developing, evolutionarily conserved unimodal regions such as primary somatomotor and visual cortex exhibit relatively strong SFC, whereas later-developing, evolutionarily expanded transmodal association regions, including limbic and default mode networks, show weaker SFC along a sensorimotor-association hierarchy\cite{marguliesSituatingDefaultmodeNetwork2016, fotiadisMyelinationExcitationinhibitionBalance2023}. Beyond these static measures, within-session variability of SFC is also heterogeneously distributed across the brain: temporal SFC variance is highest in limbic systems, intermediate in default mode, frontoparietal, dorsal and ventral attention networks, and lowest in primary visual and somatomotor cortices, reflecting how the expression of structural wiring in functional interactions can dynamically reconfigure even within a single recording session\cite{fotiadisMyelinationExcitationinhibitionBalance2023}.
Recent studies in school-age children and adolescents suggest that SFC continues to reconfigure along similar hierarchies, but these efforts have largely focused on older age ranges, coarse region-of-interest parcellations, and single-modality characterizations of SC or FC\cite{xiaDevelopmentSensorimotorvisualConnectome2024,baumDevelopmentStructureFunction2020}. In contrast, the organization and maturation of SFC during the first five years of life---particularly at vertex-level resolution and unconstrained by parcellations---remain largely unexplored.


% Recent studies in school-age children and adolescents suggest that SFC continues to reconfigure along similar hierarchies, but these efforts have largely focused on older age ranges, coarse region-of-interest parcellations, and single-modality characterizations of SC or FC\cite{xiaDevelopmentSensorimotorvisualConnectome2024,baumDevelopmentStructureFunction2020}. The organization and maturation of SFC in the first five years of life---particularly at the vertex-level resolution uncontrained by parcellations---remain essentially explored.

% The organization and maturation of SFC during the first five years of life—particularly at vertex-level resolution, without constraints from parcellations—remain largely unexplored.



% Microstructure, dynamics, and cortical gradients (myelin, E/I, Hurst, SA axis)
Intracortical myelination and excitation-inhibition (E/I) balance are key microcircuit properties that shape how structural architecture supports functional dynamics\cite{fotiadisMyelinationExcitationinhibitionBalance2023}. Myelinated axons aggregate into long-range white-matter tracts that scaffold cortical communication, and their developmental refinement—via increases in myelin and changes in axon caliber—proceeds on tract-specific timetables, with sensorimotor cortex myelinating earlier and association cortex following more protracted trajectories\cite{millerProlongedMyelinationHuman2012,deoniCorticalMaturationMyelination2015a}. Myelination thus indexes maturation of the structural substrate that constrains communication, providing a biologically grounded measurement of where the scaffold becomes efficient earliest along a sensorimotor-association (SA) axis\cite{baumGradedVariationT1w2022,yuInitialMyelinationCentral2023}. Complementing this structural view, E/I balance influences the temporal regime of local activity: the Hurst exponent captures long-range temporal correlations in BOLD fluctuations and has been proposed as a non-invasive proxy for E/I, with lower Hurst exponent values reflecting heightened E/I ratio and more desynchronized dynamics\cite{fotiadisMyelinationExcitationinhibitionBalance2023,bruiningMeasurementExcitationinhibitionRatio2020}. Computational and empirical work links Hurst exponent to inhibitory tone, although its interpretation as a direct E/I marker remains debated\cite{nishioHurstExponentMarker2024}. At the macroscale, the SA gradient describes a principal axis of cortical organization running from early-developing unimodal sensorimotor regions to later-developing transmodal association areas, along which myelination, E/I-related markers, and SFC all vary systematically in adults and older youth\cite{fotiadisMyelinationExcitationinhibitionBalance2023,larsenDevelopmentalReductionExcitation2022}. However, how these microstructural and dynamical features co-develop along the SA axis in the first years of life, and how they jointly shape the emergence of structure-function coupling, is essentially unknown.

The earliest postnatal months represent a critical window for the emergence of large-scale brain organization, yet functional gradient analysis in early infancy has been limited by the scarcity of large, high-quality datasets and the unique challenges of infant fMRI, including motion sensitivity, low signal-to-noise, and rapidly changing anatomy.

% %  Specific gap in the field & what is needed
% Despite these advances, several critical gaps remain in our understanding of how cortical structure and function become aligned in early life. First, no study has systematically characterized structure-function coupling in the first five years of life, when cortical circuits are undergoing the most rapid and plastic reconfiguration. Developmental work on SC and FC has largely focused on older children and adolescents, typically using region-of-interest parcellations and single-modality analyses that obscure fine-grained spatial heterogeneity\cite{baumDevelopmentStructureFunction2020,xuStructuralConnectivityMatures2024}. Related multimodal studies in adults and older youth have linked SFC to microstructural and dynamical features along the sensorimotor-association axis\cite{xuStructuralConnectivityMatures2024}, but they do not address how these relationships emerge during infancy. In particular, there is no vertex-wise, multimodal framework that jointly interrogates SC, FC, SFC, intracortical myelination, and E/I balance along the SA gradient in children from birth to 6 years of age. As a result, it remains unknown whether a sensorimotor-association hierarchy organizes the co-development of structural wiring, functional coupling, and microcircuit properties from birth, and how this early organization scaffolds the later functional specialization of transmodal association cortex.

%  This study: approach, key idea, and main advances
Despite these advances, key questions remain about how structure-function coupling emerges during the first five years of life. Here, we analyze multimodal MRI from birth to age 5 to derive vertex-wise SC, FC, SFC, intracortical myelination (T1w/T2w), and E/I balance (indexed by the Hurst exponent), and we integrate findings across vertices, Schaefer atlas systems, canonical networks, and the continuous sensorimotor-association axis. We address five aims. First, we characterize global developmental trajectories of SFC, myelination, and E/I balance. Second, we test whether SC, FC, and SFC follow a shared hierarchical organization across spatial scales and along the SA axis. Third, we determine how intracortical myelination relates to SFC across development. Fourth, we determine how E/I balance relates to SFC across development. Fifth, we integrate these findings by quantifying spatial and temporal SFC coupling with myelination and E/I balance, and by testing mediation models of how microstructure and E/I-related dynamics shape the expression of SC-FC co-development as SFC.
